%% kapa: exam=1 round=35 year=2024 period=2 form=A subject=appraisal_law subject_ko=감정평가관계법규 seq=9 number=9
%% source: Q-Net 감정평가사 기출문제 (cst00309) artlSeq=5213229 file=raw/1008/5213229/2교시 문제지 원본.pdf
%% conversion: kordoc 4.24.51 -> tools/gen1.py (auto); review=proofread-sonnet; flags=none
\begin{question}{9}
국토의 계획 및 이용에 관한 법령상 개발밀도관리구역에 관한 설명으로 옳지 않은 것은?
\begin{choices}
  \item 공업지역에서의 개발행위로 기반시설의 수용능력이 부족할 것이 예상되는 지역 중 기반시설의 설치가 곤란한 지역을 개발밀도관리구역으로 지정할 수 있다.
  \item 개발밀도관리구역에서는 해당 용도지역에 적용되는 용적률 최대한도의 30퍼센트 범위에서 용적률을 강화하여 적용한다.
  \item 개발밀도관리구역을 변경하려면 해당 지방자치단체에 설치된 지방도시계획위원회의 심의를 거쳐야 한다.
  \item 지정권자는 개발밀도관리구역을 지정한 경우 그 사실을 당해 지방자치단체의 공보에 게재하는 방법으로 고시하여야 한다.
  \item 개발밀도관리구역의 지정기준을 정할 때 고려되는 기반시설에 수도공급설비도 포함된다.
\end{choices}
\end{question}
\answer{2}
